% GENERATED FILE. Edit canonical lessons, then run npm run generate:books.
% canonical-source-hash: fnv1a64:f5336020537b3ba2
% canonical-lessons: KA-C310-svaccha, KA-S174-letter-chha, KA-C310-adara-nantara, KA-C310-bandu, KA-C310-atithigalu, KA-R310-a-story-in-order

\chapter{Telling a Story in Order}
\label{ch:ka-telling-a-story-in-order}
\hlchaptermodality{310}

\begin{chapteropening}
\textbf{By the end of this chapter:} \emph{I can tell a short story in order with modalu, adara nantara and konege, join actions with bandu and hōgi, read a connected story of about ninety words, and read the letter \kn{ಛ} in svaccha.}

Complete the last lesson of chapter 310: tell a story in order with first, after that, having done and in the end, and read and write \kn{ಛ}.
\end{chapteropening}

\section[\texorpdfstring{\kn{ಸ್ವಚ್ಛ} (svaccha)}{ಸ್ವಚ್ಛ (svaccha)}]{\kn{ಸ್ವಚ್ಛ} (svaccha) — clean}

\label{lesson:KA-C310-svaccha}



\begin{quote}
Before the new one: say \textquotedblleft{}to sweep\textquotedblright{} and \textquotedblleft{}to wash\textquotedblright{}. (\emph{guḍisu}, \emph{toḷe}.) Then \textquotedblleft{}There was a small shop in our town.\textquotedblright{} (\emph{namma ūrinalli ondu cikka aṅgaḍi ittu}.)
\end{quote}

\subsection*{You'll want to know: \kn{ಸ್ವಚ್ಛ}}
\textbf{\kn{ಸ್ವಚ್ಛ}} — \emph{svaccha} — \textquotedblleft{}clean\textquotedblright{}.

\begin{quote}
\textbf{\kn{ನಿನ್ನೆ} \kn{ನಾನು} \kn{ನನ್ನ} \kn{ಕೋಣೆ} \kn{ಸ್ವಚ್ಛ} \kn{ಮಾಡಿದೆ}.}
\end{quote}

\begin{quote}
\emph{ninne nānu nanna kōṇe svaccha māḍide.} — \textquotedblleft{}Yesterday I cleaned my room.\textquotedblright{}
\end{quote}

\begin{quote}
\textbf{\kn{ಈಗ} \kn{ಕೋಣೆ} \kn{ಸ್ವಚ್ಛವಾಗಿದೆ}.}
\end{quote}

\begin{quote}
\emph{īga kōṇe svacchavāgide.} — \textquotedblleft{}Now the room is clean.\textquotedblright{}
\end{quote}

\textbf{\kn{ಸ್ವಚ್ಛ} \kn{ಮಾಡು}}, \textquotedblleft{}to make clean\textquotedblright{}, is how Kannada says \textquotedblleft{}to clean\textquotedblright{}, the way \textbf{\kn{ಕೆಲಸ} \kn{ಮಾಡು}} says \textquotedblleft{}to work\textquotedblright{}. To say something \textbf{is} clean, add \textbf{\kn{ಆಗಿದೆ}}, \textquotedblleft{}has become, is\textquotedblright{}, joined by a \textbf{-\kn{ವ}-}: \textbf{\kn{ಸ್ವಚ್ಛವಾಗಿದೆ}}. In front of a noun it takes \textbf{-\kn{ವಾದ}}: \textbf{\kn{ಸ್ವಚ್ಛವಾದ} \kn{ನೀರು}} (\emph{svacchavāda nīru}, \textquotedblleft{}clean water\textquotedblright{}).

Say the middle slowly: \emph{svac-cha}. Hold the \emph{c}, then let a puff of breath out with the second one. That second consonant is a letter you have not met yet, written small under \textbf{\kn{ಚ}}. The next lesson takes it out.

\begin{cousinweb}
\textbf{\kn{ಸ್ವಚ್ಛ}} (\emph{svaccha}) is Sanskrit, kept whole as a tatsama. Monier-Williams gives \emph{svaccha} as \textquotedblleft{}very transparent or clear\textquotedblright{}, and also \textquotedblleft{}pure\textquotedblright{}, and writes it \emph{sv-accha}: \emph{su-}, \textquotedblleft{}well, very\textquotedblright{}, and \emph{accha}, \textquotedblleft{}clear\textquotedblright{}. The Kannada sense, \textquotedblleft{}clean\textquotedblright{}, sits close to that \textquotedblleft{}clear, pure\textquotedblright{}: nothing in the way, nothing on it.
\end{cousinweb}

\subsection*{Your turn}
Say these aloud:

\begin{itemize}
\raggedright
  \item \emph{svaccha}, slowly, then at speed
  \item \textquotedblleft{}Yesterday I cleaned my room\textquotedblright{} — \emph{ninne nānu nanna kōṇe svaccha māḍide}
  \item \textquotedblleft{}Now the room is clean\textquotedblright{} — \emph{īga kōṇe svacchavāgide}
  \item \textquotedblleft{}clean water\textquotedblright{} — \emph{svacchavāda nīru}
\end{itemize}

\begin{tcolorbox}[breakable,colback=teal!4,colframe=teal!35!black,title={Before you move on}]
What does \textbf{\kn{ಸ್ವಚ್ಛ}} mean? (\textbf{Clean.}) \textquotedblleft{}To clean\textquotedblright{}? (\textbf{\kn{ಸ್ವಚ್ಛ} \kn{ಮಾಡು}}.) \textquotedblleft{}It is clean\textquotedblright{}? (\textbf{\kn{ಸ್ವಚ್ಛವಾಗಿದೆ}}.) Which language is it from? (\textbf{Sanskrit}: \textquotedblleft{}very clear\textquotedblright{}.)
\end{tcolorbox}
\section[\texorpdfstring{\kn{ಛ} (cha)}{ಛ (cha)}]{\kn{ಛ} (cha) — the breath partner of \kn{ಚ}, taken from \kn{ಸ್ವಚ್ಛ}}

\label{lesson:KA-S174-letter-chha}



\begin{quote}
Before the new one: which letter is \textbf{\kn{ಚ}}? (\emph{ca}.) And \textbf{\kn{ಧ}}? (\textbf{\kn{ದ}} with the breath let out, \emph{dha}.)

Now say \textbf{\kn{ಸ್ವಚ್ಛ}} (\emph{svaccha}), \textquotedblleft{}clean\textquotedblright{}. You met it a moment ago. Its second \emph{c} carries the letter of this lesson.
\end{quote}

\subsection*{Script you'll notice: \kn{ಛ}}
\textbf{\kn{ಛ}} — \emph{cha}.

It is \textbf{\kn{ಚ}} (\emph{ca}) with a puff of breath after it. Hold a hand in front of your mouth: on \emph{ca} the air stays put, on \emph{cha} it comes out. The \emph{ch} of the romanization is that breath. It is not the English \emph{ch} of \emph{church}; that sound is already \textbf{\kn{ಚ}}.

It is the pair the breath series was missing. You can read the others, from \textbf{\kn{ಕ} \kn{ಖ}} and \textbf{\kn{ಗ} \kn{ಘ}} to \textbf{\kn{ಪ} \kn{ಫ}} and \textbf{\kn{ಬ} \kn{ಭ}}; now \textbf{\kn{ಚ} \kn{ಛ}} joins them.

Its shape does not grow out of \textbf{\kn{ಚ}} the way \textbf{\kn{ಧ}} grows out of \textbf{\kn{ದ}}. In print \textbf{\kn{ಛ}} has a small closed loop on the left, a round, cup-shaped body, a right side that rises and bends over to the right at the top, and a short straight tail hanging from the bottom of the cup. That tail is the breath mark you know from \textbf{\kn{ಧ}} and \textbf{\kn{ಢ}}.

In \textbf{\kn{ಸ್ವಚ್ಛ}} you met it small: \textbf{\kn{ಚ್ಛ}} is \textbf{\kn{ಚ}} with a reduced \textbf{\kn{ಛ}}, without its top hook, written under it. That is how this script writes two consonants with no vowel between them, as \textbf{\kn{ವ}} sits small under \textbf{\kn{ಸ}} in \textbf{\kn{ಸ್ವ}}. Most words that hold \textbf{\kn{ಛ}} came from Sanskrit.

\subsection*{Writing: \kn{ಛ} — copy the printed shape}
Put your pen beside \textbf{\kn{ಛ}} and copy it, slowly, and larger than it is printed. Close the small loop on the left before you go round the cup, and keep the tail short and straight.

\begin{quote}
This book does not yet tell you where to start this character or which way to travel. That is a real thing, taught with real variation from school to school, and it is not written down here until it can be written down with a source. Copying the printed shape needs no such source, and it is how the shape becomes yours in the meantime.
\end{quote}

\subsection*{Your turn}
\begin{itemize}
\raggedright
  \item \emph{Find:} \kn{ಛ}, small, under \kn{ಚ} in \textbf{\kn{ಸ್ವಚ್ಛ}}
  \item \emph{Trace it:} \kn{ಛ} three times, saying \emph{cha} as you finish each one
  \item \emph{Write it:} \kn{ಚ}, then \kn{ಛ}, from memory, naming each
  \item \emph{Write it:} \textbf{\kn{ಸ್ವಚ್ಛ}}
  \item \emph{Read it:} \textbf{\kn{ಸ್ವಚ್ಛ}} aloud once more
\end{itemize}

\begin{tcolorbox}[breakable,colback=teal!4,colframe=teal!35!black,title={Before you move on}]
Which letter is this — \kn{ಛ}? (\emph{\textbf{cha}}: \textbf{\kn{ಚ}} with the breath.) Which word did it come from? (\textbf{\kn{ಸ್ವಚ್ಛ}}, \emph{svaccha}, \textquotedblleft{}clean\textquotedblright{}.) Where does it sit in that word? (\textbf{Small, under \kn{ಚ}.})
\end{tcolorbox}
\section[\texorpdfstring{\kn{ಅದರ} \kn{ನಂತರ} (adara nantara)}{ಅದರ ನಂತರ (adara nantara)}]{\kn{ಅದರ} \kn{ನಂತರ} (adara nantara) — first, after that, in the end}

\label{lesson:KA-C310-adara-nantara}



\begin{quote}
Say \textquotedblleft{}afterwards\textquotedblright{} (\emph{āmēle}) and \textquotedblleft{}later\textquotedblright{} (\emph{nantara}). Then the word for clean. (\emph{svaccha}.) Which letter sits small inside it? (\textbf{\kn{ಛ}}, \emph{cha}.)
\end{quote}

\subsection*{You'll want to know: \kn{ಮೊದಲು} … \kn{ಅದರ} \kn{ನಂತರ} … \kn{ಕೊನೆಗೆ}}
\begin{quote}
\textbf{\kn{ಮೊದಲು} \kn{ನಾನು} \kn{ಕಾಫಿ} \kn{ಕುಡಿದೆ}. \kn{ಅದರ} \kn{ನಂತರ} \kn{ಶಾಲೆಗೆ} \kn{ಹೋದೆ}.}
\end{quote}

\begin{quote}
\emph{modalu nānu kāphi kuḍide. adara nantara śālege hōde.} — \textquotedblleft{}First I drank coffee. After that I went to school.\textquotedblright{}
\end{quote}

\textbf{\kn{ಮೊದಲು}} (\emph{modalu}, \textquotedblleft{}first, at first\textquotedblright{}) opens a story. \textbf{\kn{ಅದರ} \kn{ನಂತರ}} is \textbf{\kn{ನಂತರ}}, \textquotedblleft{}later\textquotedblright{}, with \textbf{\kn{ಅದರ}} (\emph{adara}, \textquotedblleft{}of that\textquotedblright{}) in front: \textquotedblleft{}after that\textquotedblright{}. \textbf{\kn{ಆಮೇಲೆ}} does the same job, and you will hear both.

\begin{quote}
\textbf{\kn{ಕೊನೆಗೆ} \kn{ರಾತ್ರಿ} \kn{ಹತ್ತು} \kn{ಗಂಟೆಗೆ} \kn{ಮನೆಗೆ} \kn{ಬಂದೆ}.}
\end{quote}

\begin{quote}
\emph{konege rātri hattu gaṇṭege manege bande.} — \textquotedblleft{}In the end I came home at ten at night.\textquotedblright{}
\end{quote}

A told story often runs on three posts: \textbf{\kn{ಮೊದಲು}} … \textbf{\kn{ಅದರ} \kn{ನಂತರ}} (or \textbf{\kn{ಆಮೇಲೆ}}) … \textbf{\kn{ಕೊನೆಗೆ}}.

\subsection*{You'll want to know: before and after a thing}
\begin{quote}
\textbf{\kn{ಊಟದ} \kn{ನಂತರ} \kn{ನಾವು} \kn{ಸಿನಿಮಾಗೆ} \kn{ಹೋದೆವು}.}
\end{quote}

\begin{quote}
\emph{ūṭada nantara nāvu sinimāge hōdevu.} — \textquotedblleft{}After the meal we went to the cinema.\textquotedblright{}
\end{quote}

\textbf{\kn{ಊಟದ}} is \textbf{\kn{ಊಟ}} with \textbf{-\kn{ದ}}, \textquotedblleft{}of\textquotedblright{}: \textquotedblleft{}after (of) the meal\textquotedblright{}. \textbf{\kn{ಮೊದಲು}} works for \textquotedblleft{}before\textquotedblright{}, but it takes the \textquotedblleft{}to\textquotedblright{} ending instead: \textbf{\kn{ಊಟಕ್ಕೆ} \kn{ಮೊದಲು}} (\emph{ūṭakke modalu}, \textquotedblleft{}before the meal\textquotedblright{}). Each word keeps its own ending.

With a number, \textbf{\kn{ನಂತರ}} counts time:

\begin{quote}
\textbf{\kn{ಎರಡು} \kn{ದಿನಗಳ} \kn{ನಂತರ} \kn{ಅವನು} \kn{ಬಂದನು}.}
\end{quote}

\begin{quote}
\emph{eraḍu dinagaḷa nantara avanu bandanu.} — \textquotedblleft{}Two days later he came.\textquotedblright{}
\end{quote}

\textbf{\kn{ದಿನಗಳ}} (\emph{dinagaḷa}, \textquotedblleft{}of days\textquotedblright{}) is \textbf{\kn{ದಿನ}} with the plural \textbf{-\kn{ಗಳ}}.

\subsection*{Your turn}
Say these aloud:

\begin{itemize}
\raggedright
  \item \textquotedblleft{}First I drank coffee. After that I went to school.\textquotedblright{} — \emph{modalu nānu kāphi kuḍide. adara nantara śālege hōde.}
  \item \textquotedblleft{}After the meal we went to the cinema\textquotedblright{} — \emph{ūṭada nantara nāvu sinimāge hōdevu}
  \item \textquotedblleft{}Two days later he came\textquotedblright{} — \emph{eraḍu dinagaḷa nantara avanu bandanu}
  \item yesterday in three steps: first, after that, in the end
\end{itemize}

\begin{tcolorbox}[breakable,colback=teal!4,colframe=teal!35!black,title={Before you move on}]
\textquotedblleft{}First\textquotedblright{}? (\textbf{\kn{ಮೊದಲು}}.) \textquotedblleft{}After that\textquotedblright{}? (\textbf{\kn{ಅದರ} \kn{ನಂತರ}}, or \textbf{\kn{ಆಮೇಲೆ}}.) \textquotedblleft{}In the end\textquotedblright{}? (\textbf{\kn{ಕೊನೆಗೆ}}.) \textquotedblleft{}Before the meal\textquotedblright{}? (\textbf{\kn{ಊಟಕ್ಕೆ} \kn{ಮೊದಲು}}.) \textquotedblleft{}After the meal\textquotedblright{}? (\textbf{\kn{ಊಟದ} \kn{ನಂತರ}}.)
\end{tcolorbox}
\section[\texorpdfstring{\kn{ಬಂದು} (bandu)}{ಬಂದು (bandu)}]{\kn{ಬಂದು} (bandu) — having come, and then}

\label{lesson:KA-C310-bandu}



\begin{quote}
Say \textquotedblleft{}First I drank coffee. After that I went to school.\textquotedblright{} (\emph{modalu nānu kāphi kuḍide. adara nantara śālege hōde.})

Two sentences, two full verbs. Kannada can fold them into one.
\end{quote}

\subsection*{You'll want to know: the -\kn{ಇ} and -\kn{ಉ} form}
\begin{quote}
\textbf{\kn{ನಾನು} \kn{ಮನೆಗೆ} \kn{ಬಂದು} \kn{ಊಟ} \kn{ಮಾಡಿದೆ}.}
\end{quote}

\begin{quote}
\emph{nānu manege bandu ūṭa māḍide.} — \textquotedblleft{}I came home and had a meal.\textquotedblright{}
\end{quote}

The first verb carries no person and no time: \textbf{\kn{ಬಂದು}}, \textquotedblleft{}having come\textquotedblright{}. Only the last verb says who and when. So a chain of these runs on to one full verb at the end:

\begin{quote}
\textbf{\kn{ಅವನು} \kn{ಅಂಗಡಿಗೆ} \kn{ಹೋಗಿ} \kn{ಒಂದು} \kn{ಪುಸ್ತಕ} \kn{ಖರೀದಿಸಿ} \kn{ಮನೆಗೆ} \kn{ಬಂದನು}.}
\end{quote}

\begin{quote}
\emph{avanu aṅgaḍige hōgi ondu pustaka kharīdisi manege bandanu.} — \textquotedblleft{}He went to the shop, bought a book and came home.\textquotedblright{}
\end{quote}

Build it from the past without its ending. Where the \textbf{-\kn{ಇ}-} came in front of the bead (\textbf{\kn{ಮಾಡಿದ}-}, \textbf{\kn{ಖರೀದಿಸಿದ}-}), keep the \textbf{-\kn{ಇ}} and drop the bead: \textbf{\kn{ಮಾಡಿ}}, \textbf{\kn{ಖರೀದಿಸಿ}}, \textbf{\kn{ನೋಡಿ}}. Everywhere else, the last \textbf{-\kn{ಅ}} becomes \textbf{-\kn{ಉ}}: \textbf{\kn{ಬಂದ}-} $\to$ \textbf{\kn{ಬಂದು}}, \textbf{\kn{ಕುಡಿದ}-} $\to$ \textbf{\kn{ಕುಡಿದು}}, \textbf{\kn{ಕೊಟ್ಟ}-} $\to$ \textbf{\kn{ಕೊಟ್ಟು}}, \textbf{\kn{ಇದ್ದ}-} $\to$ \textbf{\kn{ಇದ್ದು}}. Two common verbs go their own way: \textbf{\kn{ಹೋಗು}} $\to$ \textbf{\kn{ಹೋಗಿ}}, \textbf{\kn{ಆಗು}} $\to$ \textbf{\kn{ಆಗಿ}}.

You have met this form inside words already: \textbf{\kn{ಕಳೆದುಹೋಗು}}, \textquotedblleft{}to get lost\textquotedblright{}, starts with \textbf{\kn{ಕಳೆದು}}, \textquotedblleft{}having been lost\textquotedblright{}, and \textbf{\kn{ದಯವಿಟ್ಟು}}, \textquotedblleft{}please\textquotedblright{}, is \textquotedblleft{}having placed compassion\textquotedblright{}.

\subsection*{You'll want to know: -\kn{ಇದ} \kn{ನಂತರ}, after doing}
\begin{quote}
\textbf{\kn{ಊಟ} \kn{ಮಾಡಿದ} \kn{ನಂತರ} \kn{ನಾನು} \kn{ಮಲಗಿದೆ}.}
\end{quote}

\begin{quote}
\emph{ūṭa māḍida nantara nānu malagide.} — \textquotedblleft{}After having the meal, I lay down.\textquotedblright{}
\end{quote}

\textbf{\kn{ಮಾಡಿದ}} is the past without its ending, \textquotedblleft{}that was done\textquotedblright{}, the same piece as the \textbf{\kn{ಬಂದ}} inside \textbf{\kn{ಬಂದಾಗ}}. Before \textbf{\kn{ನಂತರ}} it gives \textquotedblleft{}after doing\textquotedblright{}: \textbf{\kn{ಬಂದ} \kn{ನಂತರ}}, \textquotedblleft{}after coming\textquotedblright{}; \textbf{\kn{ನೋಡಿದ} \kn{ನಂತರ}}, \textquotedblleft{}after seeing\textquotedblright{}. Where \textbf{-\kn{ಇ}} and \textbf{-\kn{ಉ}} tie two actions tightly, \textbf{-\kn{ಇದ} \kn{ನಂತರ}} marks a clear step in time.

\subsection*{Your turn}
Say these aloud:

\begin{itemize}
\raggedright
  \item \textquotedblleft{}I came home and had a meal\textquotedblright{} — \emph{nānu manege bandu ūṭa māḍide}
  \item \textquotedblleft{}He went to the shop, bought a book and came home\textquotedblright{} — \emph{avanu aṅgaḍige hōgi ondu pustaka kharīdisi manege bandanu}
  \item \textquotedblleft{}I drank coffee and went to school\textquotedblright{} — \emph{nānu kāphi kuḍidu śālege hōde}
  \item \textquotedblleft{}After having the meal, I lay down\textquotedblright{} — \emph{ūṭa māḍida nantara nānu malagide}
\end{itemize}

\begin{tcolorbox}[breakable,colback=teal!4,colframe=teal!35!black,title={Before you move on}]
\textquotedblleft{}Having come\textquotedblright{}? (\textbf{\kn{ಬಂದು}}.) \textquotedblleft{}Having done\textquotedblright{}? (\textbf{\kn{ಮಾಡಿ}}.) \textquotedblleft{}Having gone\textquotedblright{}? (\textbf{\kn{ಹೋಗಿ}}.) Which verb in the chain says who and when? (\textbf{Only the last.}) \textquotedblleft{}After seeing\textquotedblright{}? (\textbf{\kn{ನೋಡಿದ} \kn{ನಂತರ}}.)
\end{tcolorbox}
\section[(reading)]{(reading) — the day the guests came}

\label{lesson:KA-C310-atithigalu}



\begin{quote}
About ninety words, nearly all met before; the few new forms are glossed after it. The writer is a young woman, telling a friend about last Sunday.
\end{quote}

\begin{tcolorbox}[breakable,skin=enhanced,colback=orange!6,colframe=orange!45!black,boxrule=0.5pt,arc=1mm,left=6pt,right=6pt,top=4pt,bottom=4pt,fonttitle=\bfseries,title={Read this}]
\begin{quote}
\kn{ಕಳೆದ} \kn{ಭಾನುವಾರ} \kn{ನಮ್ಮ} \kn{ಮನೆಗೆ} \kn{ಅತಿಥಿಗಳು} \kn{ಬಂದರು}. \kn{ಮೊದಲು} \kn{ನಾನು} \kn{ಮತ್ತು} \kn{ನನ್ನ} \kn{ತಂಗಿ} \kn{ಮನೆ} \kn{ಸ್ವಚ್ಛ} \kn{ಮಾಡಿದೆವು}. \kn{ನಾನು} \kn{ಕೋಣೆ} \kn{ಗುಡಿಸಿದೆ}. \kn{ತಂಗಿ} \kn{ಪಾತ್ರೆ} \kn{ತೊಳೆದಳು}. \kn{ಅಮ್ಮ} \kn{ಅಡುಗೆ} \kn{ಮಾಡುತ್ತಿದ್ದರು}.
\end{quote}

\begin{quote}
\kn{ನಾವು} \kn{ಕೆಲಸ} \kn{ಮಾಡುವಾಗ} \kn{ಫೋನ್} \kn{ಬಂತು}. \kn{ಅಜ್ಜಿ} \kn{ಮತ್ತು} \kn{ಅಜ್ಜ} \kn{ಬಸ್} \kn{ನಿಲ್ದಾಣದಲ್ಲಿ} \kn{ಇದ್ದರು}. \kn{ಅಪ್ಪ} \kn{ನಿಲ್ದಾಣಕ್ಕೆ} \kn{ಹೋಗಿ} \kn{ಅವರನ್ನು} \kn{ಕರೆದುಕೊಂಡು} \kn{ಬಂದರು}. \kn{ನಾವು} \kn{ಒಟ್ಟಿಗೆ} \kn{ಊಟ} \kn{ಮಾಡಿದೆವು}.
\end{quote}

\begin{quote}
\kn{ಊಟದ} \kn{ನಂತರ} \kn{ಅಜ್ಜ} \kn{ಕಾಫಿ} \kn{ಕುಡಿದು} \kn{ಪತ್ರಿಕೆ} \kn{ಓದಿದರು}. \kn{ಅಜ್ಜಿ} \kn{ಒಂದು} \kn{ಕಥೆ} \kn{ಹೇಳಿದರು}. \kn{ಅವರ} \kn{ಹಳ್ಳಿಯಲ್ಲಿ} \kn{ಒಂದು} \kn{ಚಿಕ್ಕ} \kn{ಶಾಲೆ} \kn{ಇತ್ತು}. \kn{ಅಜ್ಜಿ} \kn{ಪ್ರತಿದಿನ} \kn{ನಡೆದು} \kn{ಶಾಲೆಗೆ} \kn{ಹೋಗುತ್ತಿದ್ದರು}. \kn{ಈಗ} \kn{ಅಲ್ಲಿ} \kn{ಒಂದು} \kn{ದೊಡ್ಡ} \kn{ಶಾಲೆ} \kn{ಇದೆ}.
\end{quote}

\begin{quote}
\kn{ಕೊನೆಗೆ} \kn{ರಾತ್ರಿ} \kn{ಹತ್ತು} \kn{ಗಂಟೆಗೆ} \kn{ಅತಿಥಿಗಳು} \kn{ಹೊರಟರು}. \kn{ಆಮೇಲೆ} \kn{ನಾನು} \kn{ಮಲಗಿದೆ}. \kn{ಅದು} \kn{ಒಳ್ಳೆಯ} \kn{ದಿನ}.
\end{quote}

Read it once without stopping, then once more, sorting every verb: event, scene, or habit.
\end{tcolorbox}

\begin{culture}
New forms: \textbf{\kn{ಅತಿಥಿಗಳು}} (\emph{atithigaḷu}, \textquotedblleft{}guests\textquotedblright{}), \textbf{\kn{ಅತಿಥಿ}} with the plural \textbf{-\kn{ಗಳು}}; \textbf{\kn{ನಿಲ್ದಾಣದಲ್ಲಿ}} (\emph{nildāṇadalli}, \textquotedblleft{}at the station\textquotedblright{}); \textbf{\kn{ಅವರನ್ನು}} (\emph{avarannu}, \textquotedblleft{}them\textquotedblright{}); \textbf{\kn{ಅವರ}} (\emph{avara}, \textquotedblleft{}her\textquotedblright{}), said with respect; \textbf{\kn{ಕರೆದುಕೊಂಡು} \kn{ಬಂದರು}} (\emph{karedukoṇḍu bandaru}, \textquotedblleft{}brought (them)\textquotedblright{}), literally \textquotedblleft{}having called (them) along, came\textquotedblright{}; \textbf{\kn{ಹೊರಟರು}} (\emph{horaṭaru}, \textquotedblleft{}set off, left\textquotedblright{}), the past of \textbf{\kn{ಹೊರಡು}}.

The scene is \textbf{\kn{ಮಾಡುವಾಗ}} and \textbf{\kn{ಅಡುಗೆ} \kn{ಮಾಡುತ್ತಿದ್ದರು}}; the habit is \textbf{\kn{ಹೋಗುತ್ತಿದ್ದರು}}, with \textbf{\kn{ಪ್ರತಿದಿನ}}; the background is \textbf{\kn{ಇತ್ತು}}, and \textbf{\kn{ಇದ್ದರು}} says where the grandparents were. The \textbf{\kn{ಇದೆ}} of \textquotedblleft{}now\textquotedblright{} is the present. Every other verb is an event. \textbf{\kn{ಹೋಗಿ}}, \textbf{\kn{ಕುಡಿದು}} and \textbf{\kn{ನಡೆದು}} carry no time of their own: the verb after them does.
\end{culture}

\subsection*{Your turn}
Say these aloud:

\begin{itemize}
\raggedright
  \item what the two sisters did first
  \item what Mother was doing, and what came while they were working
  \item what Grandmother used to do every day
\end{itemize}

\begin{tcolorbox}[breakable,colback=teal!4,colframe=teal!35!black,title={Before you move on}]
Why is \textbf{\kn{ಇತ್ತು}} in the story? (\textbf{The small school was there then: it is background.}) Which word closes the story's three posts? (\textbf{\kn{ಕೊನೆಗೆ}}.)
\end{tcolorbox}
\section[(review)]{Review — a story in order}

\label{lesson:KA-R310-a-story-in-order}



\begin{quote}
Say \textquotedblleft{}clean\textquotedblright{}. (\emph{svaccha}.) Which letter sits small inside it, and what does it add to \textbf{\kn{ಚ}}? (\textbf{\kn{ಛ}}: the breath.) Then \textquotedblleft{}At nine last night I was reading a book.\textquotedblright{}

\begin{itemize}
\raggedright
  \item \emph{Write it:} \textbf{\kn{ಛ}} and \textbf{\kn{ಸ್ವಚ್ಛ}}, once each from memory
\end{itemize}
\end{quote}

\subsection*{Your turn}
Say these aloud:

\begin{itemize}
\raggedright
  \item First we cleaned the house. After that Mother cooked.
  \item Father went to the station and brought Grandmother.
  \item After the meal, Grandmother told a story.
  \item In the end, at ten at night, the guests left.
  \item the day the guests came, in your own words, in four sentences
\end{itemize}

\begin{tcolorbox}[breakable,colback=teal!4,colframe=teal!35!black,title={Before you move on}]
\textquotedblleft{}Clean\textquotedblright{}? (\textbf{\kn{ಸ್ವಚ್ಛ}}.) \textquotedblleft{}First … after that … in the end\textquotedblright{}? (\textbf{\kn{ಮೊದಲು}} … \textbf{\kn{ಅದರ} \kn{ನಂತರ}} … \textbf{\kn{ಕೊನೆಗೆ}}.) \textquotedblleft{}Having come\textquotedblright{}? (\textbf{\kn{ಬಂದು}}.) \textquotedblleft{}After eating\textquotedblright{}? (\textbf{\kn{ಊಟ} \kn{ಮಾಡಿದ} \kn{ನಂತರ}}.)
\end{tcolorbox}
